\section{气体压强的微观解释}
\subsection{从分子运动到器壁受力}
宏观气体看起来连续，微观上却由大量分子组成。气体压强不是单个分子的连续推力，而是许多分子与器壁碰撞所传递动量的统计结果。单次碰撞的力随时间剧烈变化，但对大量分子、足够长时间取平均后，压强可以稳定地测量。

本讲模型把分子视为具有质量的小粒子，除碰撞瞬间外忽略分子间作用；容器中分子数很多，平衡时各方向地位相同。理想气体近似并不表示分子不碰撞，而是忽略分子间相互作用对状态方程和势能的影响。

估计分子间平均距离可用 $d\sim(V/N)^{1/3}$。该式相当于把每个分子平均分配到一个体积为 $V/N$ 的小格中，只给数量级，不意味着分子排列成晶格。若这一距离不再明显大于分子尺度，或气体接近凝结，模型的适用性就应重新检查。


取边长 $L$ 的立方体容器。一个质量为 $m$ 的分子速度分量为 $v_x,v_y,v_z$，与垂直 $x$ 轴的壁面弹性碰撞后，$v_x$ 反向，其余分量不变。器壁每次获得的动量为 $2m|v_x|$，相邻两次撞击同一壁的时间为 $2L/|v_x|$。

\fig{fig:micro}{只画分子运动在$x$方向的投影：撞右壁后反向，两次撞同一壁须走过$2L$}{\begin{tikzpicture}

\draw[world](0,-.3)--(0,2.8) (4,-.3)--(4,2.8);\fill(3,2)circle(.07);\draw[flow](3,2)--(3.9,2)node[midway,above]{$+|v_x|$};\fill(3,.7)circle(.07);\draw[flow](3,.7)--(2.1,.7)node[midway,above]{$-|v_x|$};\node[left]at(0,2){碰前};\node[left]at(0,.7){碰后};\draw[dim](0,-.8)--(4,-.8)node[midway,below]{$L$};
\node[align=left]at(7.1,1.5){分子：$\Delta p_x=-2m|v_x|$\\器壁：$\Delta P_x=+2m|v_x|$\\往返：$\Delta t=2L/|v_x|$};

\end{tikzpicture}}

图\ref{fig:micro}中，一个分子对该壁的平均力为 $mv_x^2/L$。把 $N$ 个分子的贡献相加并除以壁面积 $L^2$，得到

\[p=\frac{m}{L^3}\sum_{j=1}^{N}v_{xj}^2=\frac{Nm}{V}\langle v_x^2\rangle.
\]
平衡、各向同性条件下，$\langle v_x^2\rangle=\langle v_y^2\rangle=\langle v_z^2\rangle$，而 $\langle v^2\rangle=\langle v_x^2+v_y^2+v_z^2\rangle$，所以

\begin{equation}\label{eq:t2pressure}
\boxed{p=\frac13n_Nm\langle v^2\rangle=\frac13\rho\langle v^2\rangle},\qquad n_N=\frac NV.
\end{equation}
这里 $n_N$ 为数密度，$\rho$ 为质量密度。因分子速率分布不同，不能把 $\langle v^2\rangle$ 直接写成 $\langle v\rangle^2$；若把所有分子都设成同一速率，也会丢失统计量的本义。图中的 $\Delta P_x$ 表示动量变化，与压强 $p$ 区分。

压强公式中的三个因素各有明确作用。保持分子运动状况不变，数密度加倍，单位时间撞壁的分子数加倍，压强就加倍；保持数密度不变，较大的均方速度既增加每次传递的动量，也增加撞击频率，最终使压强正比于速度的平方。三分之一则来自三维空间中动能在三个方向的平均分配，不能在气体存在显著定向流动时未经检查地沿用。

这个推导采用简化的往返模型，但核心在于动量通量。分子间碰撞会改变具体轨迹，在平衡气体中并不改变最终的统计压强公式；宏观公式不依赖每个分子始终沿固定路线往返。

\section{温度、平动动能与理想气体内能}
\subsection{温度决定怎样的分子能量}
把式\eqref{eq:t2pressure}与 $pV=NkT$ 比较，得到
\[\frac12m\langle v^2\rangle=\frac32kT.\]
左侧是单个分子的平均平动动能。因此同温的不同理想气体具有相同的平均平动动能，质量较小的分子则需要更大的均方速度。方均根速率为

\[\boxed{v_{\mathrm{rms}}=\sqrt{\langle v^2\rangle}=\sqrt{\frac{3kT}{m}}=\sqrt{\frac{3RT}{\mu}}}.\]
升温时增加的是平均能量与均方速度，不表示每个分子都加速。某一时刻仍有快慢不同的分子；不同种类分子同温，也不意味着它们逐个配对后动能都相同。

“平均速度为零”与“平均动能为零”是完全不同的说法。平衡时向左、向右的速度分量在有符号平均中抵消，因而 $\langle\bm v\rangle=0$；但平方后各项均为非负，$\langle v^2\rangle$ 仍大于零。温度联系的是后一种统计量。因此静止容器中的气体可以没有整体动量，却具有很大的内能；气体吹向一侧时，还应把整体流速与相对整体运动的无规则速度区分。

应区分“平动动能”与“总内能”。双原子分子除质心平动外，还可能转动、振动；温度给定时，其平动平均能量仍为 $3kT/2$，但总内能一般高于只含平动的单原子气体。

\subsection{整体流动与无规则运动的分离}
压强、温度与分子速度的关系中，速度应相对气体的局部平均运动来定义。若气体具有整体流速 $\bm U$，把第 $j$ 个分子的实验室速度分解为
\[
 \bm v_j=\bm U+\bm w_j,\qquad
 \bm U=\langle\bm v\rangle,\qquad\langle\bm w\rangle=0.
\]
图\ref{fig:random-velocity}中，$\bm U$ 是所有分子共有的整体运动部分，$\bm w_j$ 是该分子相对整体的无规则运动。单个分子的 $\bm w_j$ 不必为零，为零的是许多分子的矢量平均。
\fig{fig:random-velocity}{实验室速度由整体流速与相对整体的无规则速度组成}{
\begin{tikzpicture}[x=1.3cm,y=1.3cm]
\draw[axis] (0,0)--(3.6,0);\draw[axis] (0,0)--(0,1.8);
\draw[axis,thick,blue] (0,0)--(2,0) node[midway,below]{$\bm U$};
\draw[axis,thick] (2,0)--(3,1.2) node[midway,right]{$\bm w_j$};
\draw[axis,thick] (0,0)--(3,1.2) node[midway,above left]{$\bm v_j$};
\node[align=left] at (5.3,.7){$\langle\bm w\rangle=0$\\$\langle w^2\rangle>0$};
\end{tikzpicture}}
平方取平均后，交叉项因 $\langle\bm w\rangle=0$ 消失，得到
\[
 \langle v^2\rangle=U^2+\langle w^2\rangle,\qquad
 \sum_j\frac12mv_j^2=\frac12NmU^2+\sum_j\frac12mw_j^2.
\]
右侧第一项是气体整体的动能，第二项才是平动热内能。在局部平衡、无规则运动各向同性的条件下，热力学压强与温度满足
\[
 p=\frac13\rho\langle w^2\rangle,\qquad
 \frac12m\langle w^2\rangle=\frac32kT.
\]
因此，一团气体整体运动得快并不自动表示温度高。若要计算流动撞击器壁的作用，还须计入整体动量通量；不能把实验室的 $\langle v^2\rangle$ 全部代入热压强公式。静止密闭容器中的平衡气体 $\bm U=0$，上述关系才回到前面的常见形式。这个分解也解释了为什么充气入口的动能与焓需要分别讨论，以及为什么整体机械能不能一律并入 $nC_VT$。

\subsection{自由度、能量均分与摩尔定容热容}
在经典能量均分适用的条件下，每个以二次形式进入能量的自由度，平均贡献 $kT/2$。设有效自由度为 $f$，则定量理想气体

\[U=\frac f2NkT=\frac f2nRT,\qquad C_V=\frac f2R.\]
摩尔定容热容表示定容、只存在体积功时，每摩尔物质温度升高1K所需热量；对定量理想气体，它也是 $C_V=(1/n)\dd U/\dd T$。整团气体的热容为 $nC_V$，不能把 $C_V$ 当作与气体量无关的总热容。

常温单原子气体通常取 $f=3$，只有三维平动；刚性双原子气体取 $f=5$，包含三平动和两转动；非线性刚性分子有三平动和三转动，常用近似为 $f=6$。沿分子轴转动的自由度与振动模式是否被激发，需要考虑分子的微观结构和温区。

\begin{center}\small\begin{tabular}{cccc}\toprule 模型&有效自由度&$C_V$&$C_p$\\\midrule 单原子&3&$3R/2$&$5R/2$\\刚性双原子&5&$5R/2$&$7R/2$\\非线性刚性分子&6&$3R$&$4R$\\\bottomrule\end{tabular}\end{center}
温度升高到振动开始明显参与时，热容可能随温度增大，不能继续把固定的 $f$ 套用于全部温区。题目明确给出 $C_V$ 时应优先按给定模型计算，而不按气体名称另取一个数值。

\fig{fig:h2:diatomic}{刚性双原子模型：三维质心平动之外，还有绕两条垂直键轴的转动}{\begin{tikzpicture}

\draw[world,very thick](-1,0)--(1,0);\fill[blue](-1,0)circle(.16);\fill[blue](1,0)circle(.16);\fill(0,0)circle(.04)node[below right]{$C$};\draw[guide](-1.8,0)--(2,0)node[right]{$x$（键轴）};\draw[axis](0,0)--(0,1.8)node[above]{$y$};\draw[axis](0,0)--(-1.2,-1.2)node[below]{$z$};
\draw[flow,red](.22,.3)--(.22,1.4)node[right]{$\omega_y$};\draw[flow,red](-.3,-.05)--(-1.1,-.85)node[left]{$\omega_z$};\draw[dim](-1,-1.7)--(1,-1.7)node[midway,below]{$\ell$};\draw[flow](-1,2.5)--(1,2.5)node[right]{质心平动};
\node[align=left]at(5.4,1.3){$K=K_C+\tfrac12I_y\omega_y^2+\tfrac12I_z\omega_z^2$\\五个有效二次项，各贡献$kT/2$};

\end{tikzpicture}}
图\ref{fig:h2:diatomic}把分子质心的运动和取向变化分开。两原子视为质点时，绕键轴转动不改变它们的位置，绕另外两轴则有非零转动惯量；在相应转动已被热激发的温区，得到三个平动加两个转动二次项。实际分子的量子激发条件决定哪些项能参与，不能单靠空间中画了三条轴便断定转动自由度必为三。


若键长也能振动，令伸长量为$q$，其振动能量近似为$\mu_r\dot q^2/2+\kappa q^2/2$，其中$\mu_r$为两原子的约化质量。这一个振动模式含动能和势能两个二次项；在经典充分激发时共贡献$kT$，相当于有效$f$再增加2。振动冻结时则不能计入，题设“刚性双原子”正是在排除这项键长变化。



表中的 $C_V$ 是摩尔热容，单位为 $\mathrm{J/(mol\,K)}$。若给出比定容热容 $c_V$，单位为 $\mathrm{J/(kg\,K)}$，则 $C_V=\mu c_V$，内能变化应写 $mc_V\Delta T$；把质量与摩尔热容相乘会造成量纲错误。


在本讲模型中，理想气体分子间势能忽略，同种、定量气体的内能只与温度有关：$U=U(T)$。由$\dd U=nC_V(T)\dd T$，在两个状态间积分得
\[
\Delta U=\int_{U_i}^{U_f}\dd U
=n\int_{T_i}^{T_f}C_V(T)\,\dd T.
\]
若$C_V$近似恒定，才能将它提出积分号，得到$nC_V[T]_{T_i}^{T_f}=nC_V(T_f-T_i)$；变热容时须在所给温区内对实际$C_V(T)$积分。

这里使用的 $C_V$ 并不要求实际过程始终定容。它首先是理想气体内能随温度变化的系数：两个平衡态给定后，无论实际是等压、等温还是其他过程，都可用 $\Delta U=nC_V\Delta T$ 求内能差。真正只适用于定容、无其他功过程的是 $Q=nC_V\Delta T$。把前一条正确的状态量公式误当成后一条对所有过程都成立的热量公式，是能量题中常见的混淆。

内能与整体机械运动不同。把气缸整体加速，会增加其质心动能，但不一定改变气体温度；把分子的无规则运动增强才体现为热内能增加。解题时若容器或活塞的初末宏观动能不为零，必须另列该机械能，不能并入 $nC_VT$。

若气体量发生变化，内能变化应直接按端点写
\[\Delta U=n_fC_VT_f-n_iC_VT_i.\]
只有 $n_i=n_f$ 时才可合并为 $nC_V\Delta T$。第三讲的充气问题正是利用这一区别；气体温度不变而物质的量增加时，总内能仍会增加。

\section{第一定律与系统边界}
\subsection{热量、功与能量守恒的符号}
对固定物质的量、只存在体积功的系统，取吸热为正、气体对外做功为正，第一定律为
\[\boxed{\Delta U=Q-W},\qquad Q=\Delta U+W.\]
压缩时 $W<0$，表示外界对气体做正功。气体吸热时，能量可以一部分增加内能，一部分以功的形式传给外界；气体膨胀时也可能放热，热量的正负不能单由体积增减决定。

热量和功都是过程量。它们描述跨越系统边界的能量传递，离开具体过程便不能说某个状态“具有热量”或“具有多少功”。内能则是状态量，温度和气体量相同的同种理想气体具有相同内能，与此前经过怎样的过程无关。


如果边界内只有气体，内部电阻传给气体的能量可作为热量；如果把气体和电阻一起纳入系统，则电源输入应列为电功。两种选择都可以，但同一份能量只能计一次。固定气体量、有电功输入时，可写

\[\Delta U=Q+W_{\text{电入}}-W_{\text{边界}}.
\]
若活塞推动外界弹簧或抬升重物，气体所做的功可能变成弹性势能、重力势能及对大气的功。若把气体、活塞、弹簧和地球都纳入系统，则这些量变成系统内部能量，应写整体的总能量变化，不能一面把它们当外功，一面又重复加入系统能量。

\fig{fig:boundaries}{同一弹簧活塞装置可只取气体，或把活塞和弹簧一并纳入；仍须保留大气的外部作用}{\begin{tikzpicture}

\draw[world](0,2)--(0,0)--(6.8,0) (0,2)--(6.8,2);\fill[blue!8](0,0)rectangle(3,2);\draw[very thick](3,0)--(3,2);\draw[world](6,0)--(6,2);\draw[decorate,decoration={zigzag,segment length=3pt,amplitude=2pt}](3,1)--(6,1);\node at(1.5,.8){气体$U$};\node at(4.5,1.5){弹簧$E_s$};
\draw[blue,guide,rounded corners](.12,.12)rectangle(2.88,1.88);\draw[red,guide,rounded corners](-.4,-.4)rectangle(6.4,2.4);\draw[flow](-1.4,.8)--(.7,.8)node[midway,above]{$Q$};\draw[force](4,2.9)--(3,2.9)node[midway,above]{$p_0S$};\draw[flow](3,-1)--(4,-1)node[midway,below]{$\Delta x$};\node[blue]at(1,-1.8){蓝框：气体};\node[red]at(4.8,-1.8){红框：气体、活塞与弹簧};

\end{tikzpicture}}

图\ref{fig:boundaries}中，蓝框只取气体，活塞传出的边界功包括推动大气、改变弹簧能量与活塞动能；红框则将后两项作为所含对象的能量变化。框线示意所含对象，大气始终视为外界。若只取气体，须沿实际边界计算压缩功；若把弹簧、活塞与气体合在一起，它们之间的交换便不再列为跨边界功，而由内能、弹性势能和宏观动能的变化体现。内部力对不同物体所做的功一般并不互相抵消，不能据此删去弹簧势能。选择系统的目的，是用已知的总能量变化替代难以逐时追踪的交换，同时保留仍来自外界的大气功等项。

绝热表示没有以热传递方式跨越指定边界的能量，并不排除机械功或电功。刚性表示外边界无体积变化，并不排除内部气体流动或电能输入。例\ref{t2-4}、例\ref{t2-5}和练习\ref{t2-13}分别训练这些边界选择。

\section{气体做功与状态图判读}
\subsection{由活塞位移推导体积功}
设活塞截面 $S$，缓慢移动 $\dd x$，气体体积变化 $\dd V=S\dd x$。机械平衡时气体作用力为 $pS$，于是

\[\delta W=pS\dd x=p\dd V,\qquad W=\int_{V_i}^{V_f}p\dd V.
\]
$\delta W$ 表示微小过程中的功，不是某个功状态函数的全微分。积分上限小于下限时自然给出负功；无需另外手动给压缩过程再加一次负号。

该公式的气体压强 $p$ 要在准静态边界平衡下使用。一般不可逆过程中，气体内部可能没有统一压强，实际边界功须根据外界负载计算，如 $W=\int p_{\mathrm{ext}}\dd V$，并注意所取系统是否包含活塞及其动能。自由膨胀面对真空的阻力为零，不能在初末态间随便画一条曲线后求图下面积。


在 $p$--$V$ 图中，功是曲线下的有向面积。等压过程是矩形面积；压强随体积线性变化时是梯形面积，$W=(p_i+p_f)(V_f-V_i)/2$；等温可逆过程则须积分，不能用端点平均压强代替双曲线下面积。

\fig{fig:work-paths}{相同端点的两条路径及其功之差}{\begin{tikzpicture}[x=1cm,y=1cm]
\draw[axis] (0,0)--(5.4,0) node[right]{$V$};\draw[axis] (0,0)--(0,3.5) node[above]{$p$};\node[below left] at (0,0){$O$};\fill[blue!12] (1,1)--(1,2.8)--(4,2.8)--cycle;\draw[ray] (1,1)--(1,2.8)--(4,2.8);\draw[ray] (1,1)--(4,2.8);\node[below] at (1,1){$A$};\node[above] at (1,2.8){$B$};\node[right] at (4,2.8){$C$};\node at (2.6,2.2){$W_{ABC}-W_{AC}$};
\end{tikzpicture}}

图\ref{fig:work-paths}两条路径端点相同，内能变化相同。上方路径做功更多，需吸收的热量也更多：$Q_{ABC}-Q_{AC}=W_{ABC}-W_{AC}$。用两路径围成的面积可以直接求差值，通常比把两段积分全部展开更清楚。

同一条图线还有路径与状态两层信息。点的位置给出 $pV$，从而给出定量理想气体的温度和内能；沿线的有向面积给出功。只有将这两者相加，才得到热量。例如一条向右下方的膨胀曲线有正功，但若 $pV$ 降低得足够快，内能减少的绝对值超过正功，气体仍可能放热。因此判定吸放热不能只看箭头是否向右，也不能只看压强是否降低。

闭合循环的初末状态相同，$\Delta U=0$，所以净热量等于净功。顺时针循环中膨胀段的平均压强较高，净功为正；逆时针则为负。循环面积给净功，不给总吸热量，也不能直接给热机效率；效率还需逐段确定热量的正负。

\subsection{$p$--$T$ 图中怎样判断体积}
在 $p$--$T$ 图上，$V=nRT/p$。从原点引向某状态点的射线斜率为 $p/T=nR/V$，故斜率越大、体积越小；两点处在同一条过原点射线上才代表体积相同。

\fig{fig:h2:ptrays}{同温两状态从原点引出的射线斜率不同，较高压强对应较小体积}{\begin{tikzpicture}

\draw[axis](0,0)--(5,0)node[right]{$T$};\draw[axis](0,0)--(0,3.8)node[above]{$p$};\draw[guide](0,0)--(4.4,1.76);\draw[guide](0,0)--(4.4,3.3);\fill(3,1.2)circle(.06)node[right]{$A$};\fill(3,2.25)circle(.06)node[right]{$B$};\draw[guide](3,0)--(3,2.25);\node[below]at(3,0){$T_*$};\draw[dim](3.8,1.2)--(3.8,2.25)node[midway,right]{$p_B>p_A$};\node at(2,-1){$V_B/V_A=p_A/p_B<1$};

\end{tikzpicture}}
图\ref{fig:h2:ptrays}中$A,B$的横坐标相同，先锁定温度和定量理想气体内能不变，再比较$p/T$。不能把从$A$指向$B$的竖直线斜率当成等容射线斜率；等容的射线必须从绝对温度原点出发。



若一段过程满足$p=aT+b$，则$V=nRT/(aT+b)$。求导时分子与分母均随$T$改变，商法则给
\[
\frac{\dd V}{\dd T}
=nR\frac{(aT+b)\cdot1-T\cdot a}{(aT+b)^2}
=\frac{nRb}{(aT+b)^2}.
\]
因此非零截距的符号会决定该段体积如何变化，不能只凭曲线“向右上方走”作判断。

判读时可以先看横坐标确定平均平动动能和内能的变化，再比较 $T/p$ 确定体积，最后如有必要转到 $p$--$V$ 图求功。例\ref{t2-3}与第三讲例1采用的就是这一顺序。

这里的射线斜率必须从绝对零温对应的原点量取；若横轴使用摄氏温度 $t$，则应先换成 $T=t+273.15$，不能从图中摄氏零点作射线判断体积。


一段路径即使在 $p$--$T$ 图中是直线，在 $p$--$V$ 图中也不一定仍是直线。由状态方程消去温度后，应重新得到真正的压强随体积关系，才能把图下面积作为功。不同坐标图保留的是同一过程，面积却具有不同单位，不能将两图的几何面积混用。


\section{热容与非标准过程}
\subsection{定容热容与定压热容的物理差别}
定容时体积功为零，输入热量全部增加气体内能：$Q_V=nC_V\Delta T$。定压时气体通常膨胀，$W=p\Delta V=nR\Delta T$，故

\[Q_p=nC_V\Delta T+nR\Delta T=nC_p\Delta T,\quad\boxed{C_p-C_V=R}.
\]
这个差值是每摩尔理想气体升高1K时多做的等压膨胀功，不是内能差。两个过程温升相同时内能增量相同，但定压过程需要更多热量。若温升、气体量或热容不同，则不能直接按 $C_p/C_V$ 比较总热量。

热容的下标标明过程条件。任意过程的$Q/\Delta T$一般不等于$nC_V$或$nC_p$。例如$p=kV$、$k>0$时，先对$pV$用乘积法则：$\dd(pV)=V\dd p+p\dd V$，再由$\dd p=k\dd V$，得到$nR\dd T=(kV+p)\dd V=2p\dd V$，所以$\delta W=nR\dd T/2$。代入第一定律后，在温度端点间积分，$Q=n(C_V+R/2)\int_{T_i}^{T_f}\dd T=n(C_V+R/2)[T]_{T_i}^{T_f}$，因此

\[
 \delta Q=n\left(C_V+\frac R2\right)\dd T,\qquad
 Q=n\left(C_V+\frac R2\right)\Delta T.
\]
这也给出 $W=nR\Delta T/2=\Delta(pV)/2$，与直线下的梯形面积一致。额外的 $R/2$ 来自该路径的膨胀功，介于等容的零与等压的 $R$ 之间；它不是分子新增了一种微观自由度。将路径造成的过程热容与物质的内能系数分开，才能理解非标准过程。

这解释了例\ref{t2-9}、例\ref{t2-10}中为什么出现不同于定容、定压的系数。应先由装置导出路径，再计算功，不能先选择一个熟悉的热容强套。


\fig{fv:h2:capacities}{相同温升下，定压过程还要向外推动活塞}{\begin{tikzpicture}[scale=.85]
\foreach \a in {0,5}{\begin{scope}[xshift=\a cm]\draw[thick](0,3)--(0,0)--(3,0)--(3,3);\fill[blue!9](0,0)rectangle(3,1.7);\draw[very thick](0,1.7)--(3,1.7);\end{scope}}
\draw[thick](-.2,1.7)--(.2,1.7) (2.8,1.7)--(3.2,1.7);\node at(1.5,2.5){固定活塞};
\draw[dashed](5,2.4)--(8,2.4);\draw[->,red](6.5,1.8)--(6.5,2.6);\node at(6.5,3){恒外压};
\node at(1.5,-.6){$Q_V=\Delta U$};\node at(6.5,-.6){$Q_p=\Delta U+p\Delta V$};\end{tikzpicture}}
图\ref{fv:h2:capacities}比较同样物质的量、相同温升的两份气体。左图的约束阻止边界位移，故没有体积功；右图的虚线标出受热后的活塞位置，外界获得$p\Delta V$。两份气体的内能增量相同，差别发生在跨边界的机械能传递，不能解释成定压时同一温度对应更大的分子平均平动动能。


\subsection{平衡分支与稳定性}
气态方程与受力平衡给出的代数关系，描述的是可能的平衡状态。缓慢改变温度时能否沿这些状态连续移动，还取决于平衡是否稳定。

以管口溢液为例，气柱长$\ell$、平衡压强$p$。向上移$x$后的等温气压为$p_g(x)=p\ell/(\ell+x)$。求导时$p,\ell$是原平衡态的常量，只有$x$改变，故
\[
p_g'(x)=p\ell\frac{0\cdot(\ell+x)-1\cdot1}{(\ell+x)^2}
=-\frac{p\ell}{(\ell+x)^2},\qquad p_g'(0)=-\frac p\ell.
\]
以原状态为展开点，$p_g(x)-p_g(0)=p_g'(0)x+O(x^2)$，所以一阶压降为$-px/\ell$。若液体同时溢出，负载压强下降$-\rho gx$。两者相减，净推力变化与

\[\left(\rho g-\frac p\ell\right)x
\]
同号。系数为负时上移后出现向下恢复力，系数为正时扰动继续增大。两种压降相等时，一阶恢复能力消失，应进一步检查临界状态。

稳定性检验应交代扰动的方向和温度条件。若扰动很慢，气体与恒温环境保持热平衡，气体压降用等温关系；若扰动很快且来不及换热，则应换成绝热压降。平衡曲线本身相同，回复力却可能不同。对溢出装置，还要注意溢出的液体不能自动回流，因此“出现恢复趋势”并不等于系统严格回到原来同一液体量的状态。它说明该侧不再自行扩大，适用于识别缓慢升温过程中可跟随的分支。

例\ref{t2-1}将平衡温度写成剩余汞柱高的函数。升温沿稳定支前进，到极大温度时失去稳定性；因此“最终液柱为零”的代数根不一定是升温可达到的临界温度。第四讲将用同一小扰动方法复核局部恢复性。

\subsection{从微观态理解过程方向}
一个宏观状态通常对应大量微观排列。玻尔兹曼关系 $S=k\ln\Omega$ 在等概率微观态条件下，把熵与可实现的微观态数联系起来。它说明统计规律怎样把大量分子的运动与宏观过程方向连接起来，而不是把分子想成有固定速度和固定轨迹的整齐队列。

系统熵可以因向外放热而减小；第二定律约束的是把系统和环境合起来后的总熵不减。第三讲将给出理想气体熵差的计算式，并用它检查仅满足能量守恒的末态是否可能。


若两个彼此独立的系统分别有 $\Omega_1,\Omega_2$ 种微观状态，合起来的状态数为乘积，取对数后熵就相加。这使微观状态的计数与宏观可相加的状态量相连。自发过程偏向可实现方式极多的宏观状态，描述的是大量粒子的统计规律，不是给单个分子规定一个只能向某方向运动的力。


\fig{fig:h2:freeexpansion}{刚性隔离容器中的自由膨胀：粒子数与总能量不变，可占据空间增大}{\begin{tikzpicture}

\foreach \j in {0,1}{\begin{scope}[yshift={-3.3*\j*1cm}]\draw[world](0,0)rectangle(6,2);\draw[guide](3,0)--(3,2);\node[left]at(-.3,1){\ifnum\j=0 初态\else 末态\fi};\end{scope}}
\foreach \xx/\yy in {.4/.4,.7/1.4,1.2/.8,1.8/1.6,2.2/.5,2.6/1.2}{\fill[blue](\xx,\yy)circle(.05);}\node at(4.5,1){真空};\draw[very thick](3,0)--(3,2);
\foreach \xx/\yy in {.6/-2.9,1.5/-1.8,2.2/-2.7,3.3/-2,4.5/-2.9,5.3/-1.7}{\fill[blue](\xx,\yy)circle(.05);}\node at(3,-4.1){$V\to2V,\quad N\ \text{不变},\quad U\ \text{不变}$};

\end{tikzpicture}}
图\ref{fig:h2:freeexpansion}中移去隔板后，右室没有阻力，外边界又固定、绝热，故总内能不变。若用稀薄理想气体的独立位置计数理解，每个分子可占据的空间加倍，$N$个分子的位置状态数增加约$2^N$倍；初末动量分布对应同一温度，其计数因子在比值中约去。因此$\Delta S=k\ln2^N=Nk\ln2>0$。少量圆点仅表示同一批分子，不能据图中的点数判断宏观气体的粒子总数。



例如理想气体在刚性隔离容器内向真空自由膨胀，初末总能量可以相同，但可占据的空间增大，熵增加。能量守恒只限制哪些状态具有相同能量，并不保证这些状态都能在相同条件下自发互相到达。将能量与熵的作用分开，才能解释许多温度不变却具有明显方向性的过程。

